% methods_data.tex
% Opening prose of the Methods section (no subsection header; it sits directly under \section{Methods}):
% data and vendor, universe and sectors, split basis, windows, growth rates, risk-free rate, initial wealth,
% cash, fills and trading costs. Data-preparation details (sector counts, split detection) are in appendix_b.tex.
% Sources (the code is the truth where docs disagree): scripts/00_fetch_data.jl, code/src/Data.jl,
% code/src/Costs.jl, code/src/Backtest.jl (_bt_fill!, _bt_adv), code/src/StudyBacktest.jl (run_study),
% config/study.toml, results/calibration_summary.toml, README.md.

We tested a rebalancing strategy built from three parts, a score that ranks stocks by expected growth, a
constrained utility maximization that sets how many shares of the selected stocks to hold, and a
model-predictive controller that decides when to re-plan, in a backtest on recorded prices and in a stress test
on synthetic prices. The recorded prices were the daily opens and closes (USD/share) and volumes (shares/day) of
the constituents of the S\&P~500 index and of the exchange-traded fund SPY, which served as the market and was
never held. They came from Polygon.io in three files covering 2014-01-03 to 2024-12-31, calendar 2025, and
2026-01-02 to 2026-09-04, and the trading days were those of SPY. Each name took one of the $11$ sectors of the
Global Industry Classification Standard (GICS) from a snapshot of the index constituent list. The universe kept
every ticker that appeared in all three files with a sector and had a positive open and close on every trading
% results/calibration_summary.toml (universe_size = 373; sector_counts from 15 to 60)
day from 2014-01-03 to 2026-09-04, which left $N = 373$ names with between $15$ and $60$ per sector
(Appendix~\ref{app:controller}), and we write $s(i)$ for the sector of name $i$. Because the snapshot lists current
members, the universe holds only names that were index members at the snapshot date and traded on every day of
$2014$--$2026$, each with its current sector. The prices were split-adjusted by the vendor and not adjusted for
dividends, and we used them as given, so neither the portfolios nor SPY received dividends. Each file was
% README.md (Data and its limitations); results/large_k/run.log (seven split-basis lines)
split-adjusted as of its own creation date, so for the seven names that split between the creation dates of two
files we put the earlier rows on the newest file's basis (Appendix~\ref{app:controller}).

% config/study.toml [window]; trading-day counts with code/src/Data.jl on data/ and talk/data/market_state.csv (420 rows)
The calibration window ran from 2014-01-03 to 2024-12-31 ($2{,}767$ trading days) and the trading window from
2025-01-02 to 2026-09-04 ($420$ trading days). Quantities fitted on the calibration window (the prior of the online
single-index fit, the growth moments of SPY, the volatility cap, the gain of the market state and the covariance
shrinkage) were fitted once and then held fixed. In the trading window every decision was made after the close of
day $t$ from data dated on or before day $t$, and its orders filled at the open of day $t+1$. We measured
performance with continuously compounded growth rates. For a price or wealth series $S$ observed on trading days,
the growth rate over the $k$ trading days that end on day $t$ is
\begin{equation}
  g(t;k) \;\equiv\; \frac{1}{k\,\Delta t}\,\ln\!\left(\frac{S(t)}{S(t-k)}\right),
  \label{eq:growth}
\end{equation}
% code/src/Data.jl (growth_rates, Delta t = 1/252); 419 = 420 - 1 daily steps
where $\Delta t = 1/252$~yr is the length of one trading day, so that $g$ has units of $\mathrm{yr}^{-1}$. We write
$g_i(t) = g_i(t;1)$ for the daily growth rate of name $i$ and $g_M(t)$ for that of SPY, and the realized growth rate
of a run is Eq.~\eqref{eq:growth} applied to its wealth over the $k = 419$ daily steps of the trading window.

The risk-free growth rate came from the 3-month Treasury bill rate DTB3~\cite{fredDTB3}, a bank discount rate $d$
(a fraction, on a 360-day year), converted to the continuously compounded growth rate of a $91$-day bill,
\begin{equation}
  g_f(t) \;=\; -\,\frac{365}{91}\,\ln\!\left(1 - \frac{91\,d}{360}\right),
  \label{eq:riskfree}
\end{equation}
in $\mathrm{yr}^{-1}$. Day $t$ used the last observation dated before $t$, because the rate for a day is published
% data/DTB3.csv through code/src/Data.jl (_risk_free): min 0.0358, max 0.0439 over 2025-01-02 to 2026-09-04
on the following day, and $g_f$ ranged from $0.036$ to $0.044$~$\mathrm{yr}^{-1}$ over the trading window. It set
the growth of cash in the backtest and in the controller's forecast, the excess growth of the tangency
benchmark, and the excess growth behind the Sharpe ratio and the regression on SPY.

% config/study.toml (W0 = 100000.0); code/src/Backtest.jl (_BT_HELD_W = 1e-3)
Each run started with $W_0 = 100{,}000$~USD in cash at the close of 2025-01-02, and its wealth at the close of day
$t$ was $W(t) = c(t) + \sum_i n_i(t)\,S_i(t)$, with cash $c$ (USD), fractional shares $n_i$ (negative for a short
position) and close prices $S_i$ (USD/share). Positive cash grew by the factor $e^{g_f(t)\,\Delta t}$ each day.
Orders filled at the next open, sells before buys, with all buys scaled down by one common factor when the cash
could not cover them, so cash never went negative, and a target position worth at most $0.1\%$ of wealth was set to
zero shares. A fill of $q$ shares (signed) of name $i$ at the open price $S^{\mathrm{o}}_i$ (USD/share) paid
\begin{equation}
  C(q) \;=\; \bar c\, S^{\mathrm{o}}_i\,|q| \;+\; \kappa\,\frac{|q|}{A_i}\,S^{\mathrm{o}}_i\,|q|
  \label{eq:cost}
\end{equation}
% config/study.toml [costs] half_spread_bps = 5.0, impact_kappa = 0.001; code/src/Backtest.jl (_BT_ADV_DAYS = 63)
in USD, the half spread $\bar c = 5$~bps of the traded value plus a linear impact cost~\cite{almgrenChriss2001}
with coefficient $\kappa = 0.001$ on the participation $|q|/A_i$, where $A_i$ (shares/day) is the mean daily volume
of name $i$ over the $63$ trading days that end on the decision day. There was no commission, and the same half
% config/study.toml [allocators] borrow_rate = 0.0075
spread prices the turnover limit of the allocation problem. Short positions, which only the 130/30 benchmark took,
paid a borrow fee at the rate $b = 0.0075$~$\mathrm{yr}^{-1}$ on their value at each close, charged daily, and their
sale proceeds were held as cash.

